% Imporditud: tools/import_eksperimendid.py
% Allikas: Eesti füüsikaolümpiaadi eksperimendi ülesannete
%   kogu (Rael Kalda, 2023), ülesanne Ü86, lahendus L86.
\setAuthor{EFO žürii}
\setRound{piirkonnavoor}
\setYear{2001}
\setNumber{G E2}
\setDifficulty{3}
\setTopic{geomeetriline optika}
\setVahendid{korgiga suletud katseklaas, milles on värvitu ja kollane vedelik}
\setPunktid{10}
\setAllikas{Eksperimendi kogu Ü86, lahendus lk 70}
\setLahendusseis{toores_ilma_skeemita}

\prob{Vedeliku murdumisnäitaja}
Hinnata, kas suurem murdumisnäitaja on värvitul või kollasel vedelikul. Põhjendage tulemust.

\hint

\solu
$E_2$

Lahendus: Tuleb vaadata läbi vedelike mingit eset või kujundit paberil. Teatud kauguse korral saame suurendatud kujutise. Mida suurem on kujutis, seda suurem on murdumisnäitaja, sest murdva pinna kumerus on mõlemal juhul ühesugune. Suuremale murdumisnäitajale vastab suurem murdumisnurk, sellele aga suurem kujutis. Parim seletus on joonise abil, kus on näha kujutise tekkimine läätses kahel juhul. Suuremale murdumisnäitajale vastab aga väiksem fookuskaugus.

\textit{Žürii hindamisskeem on kogumikus lk 70; siia ei ole seda ümber kirjutatud.}
\probend
